\chapter*{Introducción}
\addcontentsline{toc}{chapter}{Introducción}

En los últimos años, el avance de la inteligencia artificial ha impulsado una transformación
significativa en la forma en que se conciben y desarrollan las aplicaciones web. Este fenómeno ha
favorecido la transición desde el paradigma de la Web~2.0, centrado en la interacción y la generación
de contenido por parte de los usuarios, hacia un modelo más dinámico y adaptativo, directamente
asociado con la Web~3.0. En este nuevo contexto, las aplicaciones no solo procesan información,
sino que incorporan mecanismos que les permiten analizar su entorno, adaptarse a condiciones
cambiantes y evolucionar de manera continua.

En este nuevo escenario, la integración de mecanismos basados en inteligencia artificial surge como
una consecuencia natural de esta evolución, permitiendo incorporar capacidades de análisis,
automatización y apoyo a la toma de decisiones dentro del sistema.

Para ello, es necesario establecer bases arquitectónicas que faciliten la incorporación progresiva
de estos mecanismos sin comprometer la estabilidad del software. Esto implica diseñar aplicaciones
que puedan escalar, mantenerse y evolucionar de forma controlada. Sin embargo, la efectividad de
esta integración depende en gran medida del estilo arquitectónico, ya que estructuras altamente
acopladas pueden limitar la incorporación de nuevos componentes, especialmente en sistemas de
mayor complejidad.

Como respuesta a esta limitación, han surgido enfoques arquitectónicos orientados a la
modularidad, entre los cuales destacan las arquitecturas basadas en microservicios. Este tipo de
arquitectura permite descomponer un sistema complejo en componentes simples e independientes,
facilitando el desarrollo incremental, el despliegue continuo y la evolución controlada de sus
funcionalidades. Sin embargo, la adopción de estos enfoques también introduce nuevos desafíos,
especialmente en lo relacionado con la gestión de la complejidad operativa, la observabilidad del
sistema y el mantenimiento en entornos distribuidos.

El sistema PortalAsig, desarrollado para la gestión académica en la Facultad de Ciencias de la
Universidad Central de Venezuela, constituye un caso representativo de estas problemáticas. Desde
su creación en 2013, el sistema fue concebido con una arquitectura centralizada, acompañada de una
interfaz de usuario desarrollada en Angular. Posteriormente, el trabajo de grado de Vásquez Balza \cite{teg_rafael}
introdujo un rediseño del frontend  en Vue bajo un enfoque orientado a la mantenibilidad, incorporando
prácticas de organización y mejora de la calidad del código.

No obstante, el backend ha permanecido basado en una arquitectura monolítica con limitaciones
estructurales significativas, lo que ha incrementado su complejidad y dificultado su evolución.
El rediseño del backend hacia una arquitectura de microservicios se plantea como una necesidad
para soportar una nueva forma de gestionar el mantenimiento del sistema.

En la actualidad, el mantenimiento del sistema se realiza principalmente de forma reactiva, donde
el personal a cargo atiende fallas bajo demanda ante la aparición de errores visibles o durante
períodos específicos, como el inicio o cierre de un semestre. Las acciones correctivas son
manuales y no responden a un proceso sistemático de diagnóstico o análisis de causa raíz, lo que
dificulta la evolución controlada del sistema y la implementación de mejoras.

La integración de mecanismos asistidos por inteligencia artificial
representa una oportunidad para transformar el proceso actual de mantenimiento del software. Estas
capacidades permiten complementar el trabajo humano mediante la automatización de tareas de
análisis, diagnóstico y apoyo en la toma de decisiones, lo cual resulta especialmente relevante en
contextos donde los recursos humanos son limitados.

En consecuencia, el presente Trabajo Especial de Grado propone el diseño e implementación de un
modelo de mantenimiento asistido por inteligencia artificial para una arquitectura de
microservicios en el sistema PortalAsig. Este modelo busca integrar principios de sistemas
auto-adaptativos, mecanismos de observabilidad y técnicas de inteligencia artificial, con el
objetivo de transformar el mantenimiento del software en un proceso semi-asistido, sistemático,
eficiente y alineado con las tendencias de la cada vez más presente Web~3.0.